%%================================================
%% Appendix E
%%================================================
\chapter[ตัวอย่าง Schema, API และ MQTT Payload]{\texorpdfstring{ภาคผนวก จ\\ตัวอย่าง Schema, API และ MQTT Payload}{Appendix E: Schema API and MQTT Payload}}
\label{app:schema_payload}

ภาคผนวกนี้แสดงตัวอย่างโครงสร้างข้อมูลและ payload ที่เกี่ยวข้องกับระบบ WheelSense เพื่อประกอบคำอธิบาย server, API, database และ MQTT ในบทที่ 3 และผลการทดสอบในบทที่ 4 ตัวอย่างนี้ใช้เพื่ออธิบายรูปแบบข้อมูลระดับระบบต้นแบบ ไม่ใช่สัญญา API ฉบับสมบูรณ์สำหรับ production

\begin{table}[H]
    \centering
    \caption{กลุ่ม schema หลักของฐานข้อมูล}
    \label{tab:app_schema_groups}
    \renewcommand{\arraystretch}{1.18}
    \small
    \setlength{\tabcolsep}{3pt}
    \begin{tabular}{|L{0.26\textwidth}|L{0.34\textwidth}|L{0.30\textwidth}|}
        \hline
        \textbf{กลุ่มข้อมูล} & \textbf{ตัวอย่าง entity} & \textbf{บทบาทในระบบ} \\ \hline
        User และ workspace & users, roles, workspaces & กำหนดสิทธิ์การเข้าใช้งานและบริบทของสถานดูแล \\ \hline
        Patient และ assignment & patients, caregivers, room assignments & เชื่อมผู้พักอาศัยกับผู้ดูแล ห้อง และอุปกรณ์ \\ \hline
        Device registry & devices, device bindings, node status & ลงทะเบียน M5StickC Plus2, mobile gateway, wearable และ BLE node \\ \hline
        Telemetry และ vitals & imu\_telemetry, mobile\_device\_telemetry, vital readings & เก็บข้อมูล IMU, RSSI, HR, PPG, step count และสถานะอุปกรณ์ \\ \hline
        Localization และ alert & room prediction, alerts, workflow audit & เก็บผลตำแหน่ง เหตุการณ์แจ้งเตือน และประวัติการดำเนินงาน \\ \hline
        AI/runtime context & chat actions, audit trail, execution result & ใช้ประกอบ AI/MCP pipeline และ execution gate \\ \hline
    \end{tabular}
\end{table}

\begin{table}[H]
    \centering
    \caption{ตัวอย่าง API และ MQTT topic ที่เกี่ยวข้อง}
    \label{tab:app_api_mqtt_topics}
    \renewcommand{\arraystretch}{1.18}
    \small
    \setlength{\tabcolsep}{3pt}
    \begin{tabular}{|L{0.32\textwidth}|L{0.18\textwidth}|L{0.40\textwidth}|}
        \hline
        \textbf{เส้นทางหรือ topic} & \textbf{ชนิด} & \textbf{หน้าที่} \\ \hline
        \texttt{/api/patients} & REST API & จัดการข้อมูลผู้พักอาศัยและข้อมูลที่ dashboard ใช้แสดงผล \\ \hline
        \texttt{/api/devices} & REST API & ลงทะเบียนและตรวจสถานะอุปกรณ์ \\ \hline
        \texttt{/api/telemetry} & REST API & อ่านข้อมูล telemetry เพื่อใช้กับ dashboard และ report \\ \hline
        \texttt{/api/chat/actions} & REST API & เก็บ action proposal, confirmation และ execution result ของ AI \\ \hline
        \texttt{WheelSense/data} & MQTT & รับ telemetry จาก wheelchair sensor หรือ gateway รุ่นเดิม \\ \hline
        \texttt{WheelSense/mobile/\{device\_id\}/telemetry} & MQTT & รับข้อมูลจาก Flutter mobile gateway \\ \hline
        \texttt{WheelSense/mobile/\{device\_id\}/register} & MQTT & ลงทะเบียน mobile gateway เข้ากับ workspace \\ \hline
        \texttt{WheelSense/mobile/\{device\_id\}/walkstep} & MQTT & รับข้อมูลจำนวนก้าวเดินจากโทรศัพท์ \\ \hline
    \end{tabular}
\end{table}

\begin{lstlisting}[caption={ตัวอย่าง payload จาก M5StickC Plus2 ผ่าน mobile gateway},label={lst:app_m5_payload}]
{
  "device_id": "M5C-WS-001",
  "patient_id": "P001",
  "timestamp": "2026-05-22T10:00:00+07:00",
  "battery": { "percent": 84, "voltage": 4.02 },
  "imu": {
    "accel": { "x": 0.02, "y": -0.01, "z": 0.98 },
    "gyro": { "x": 0.12, "y": 0.04, "z": -0.02 },
    "magnitude": 0.981
  },
  "motion": {
    "state": "normal_motion",
    "velocity_mps": 0.42,
    "distance_m": 12.4
  }
}
\end{lstlisting}

\begin{lstlisting}[caption={ตัวอย่าง payload จาก Flutter mobile gateway},label={lst:app_mobile_gateway_payload}]
{
  "device_id": "MOBILE-GW-001",
  "patient_id": "P001",
  "timestamp": "2026-05-22T10:00:02+07:00",
  "m5": { "device_id": "M5C-WS-001", "battery_percent": 84 },
  "polar": { "hr_bpm": 76, "rr_ms": 812, "ppg_quality": "usable" },
  "phone": { "steps_today": 1840 },
  "rssi": [
    { "node_id": "BLE-NODE-01", "rssi": -61 },
    { "node_id": "BLE-NODE-02", "rssi": -74 }
  ]
}
\end{lstlisting}

\begin{lstlisting}[caption={ตัวอย่าง action candidate ที่ต้องผ่าน execution gate},label={lst:app_ai_action_candidate}]
{
  "intent": "care_alert_proposal",
  "patient_id": "P001",
  "summary": "Heart rate is above the configured threshold",
  "evidence": ["latest_vitals", "active_alerts", "patient_context"],
  "proposed_action": {
    "type": "notify_caregiver",
    "target_role": "head_nurse",
    "requires_confirmation": true
  },
  "policy": { "medical_claim": false, "execute_without_user": false }
}
\end{lstlisting}
